\sectionkicker{Source-backed technical notes}
\subsection*{Native Multimodalityと生成Media — 4o Image・Qwen Omni・Llama 4・Veo 3/Imagen 4: Technical Notes}
本欄は記事本文で圧縮した一次資料上の情報を、比較・再検証しやすい形で整理したものである。時系列、確認済みの主張、留意点、一次資料URLを掲載する。
\medskip
\subsection*{Theme at a glance}
\addcontentsline{toc}{subsection}{Theme at a glance}
\begin{center}
\footnotesize
\begin{tabularx}{\linewidth}{@{}p{0.20\linewidth}p{0.10\linewidth}p{0.14\linewidth}X@{}}
\toprule
資料 & 位置づけ & 種別 & 時系列 \\
\midrule
GPT-4o native image generation & 主要資料 & - & 2025-03-25   \\
Qwen2.5-Omni & 主要資料 & - & 2025-03-27   \\
Veo 3, Imagen 4, Flow, and Lyria 2 & 主要資料 & - & 2025-05-20   \\
\bottomrule
\end{tabularx}
\end{center}
\normalsize

\begin{technicalnote}{GPT-4o native image generation}{主要資料}
\begingroup
% reader-facing Technical Notes break policy
\widowpenalty=10000
\clubpenalty=10000
\displaywidowpenalty=10000
\begin{tabularx}{\linewidth}{@{}>{\bfseries}p{0.22\linewidth}X@{}}
組織 & OpenAI \\
種別 & - \\
時系列 & 2025-03-25 \\
\end{tabularx}
\smallskip
{\bfseries 一次資料から整理したtechnical points}
\begin{itemize}[leftmargin=1.5em,itemsep=0.35em]
\item \textbf{一次情報で確認できる事実}: 2025年3月25日、OpenAIはGPT-4oにnative image generationを統合した。
\end{itemize}
\begin{samepage}
{\bfseries 一次資料}
\begingroup\sloppy
\Urlmuskip=0mu plus 2mu\relax
\begin{itemize}[leftmargin=1.5em,itemsep=0.25em]
\item {\scriptsize\url{https://openai.com/index/introducing-4o-image-generation/}}
\end{itemize}
\endgroup
\end{samepage}
\endgroup
\end{technicalnote}
\medskip

\begin{technicalnote}{Qwen2.5-Omni}{主要資料}
\begingroup
% reader-facing Technical Notes break policy
\widowpenalty=10000
\clubpenalty=10000
\displaywidowpenalty=10000
\begin{tabularx}{\linewidth}{@{}>{\bfseries}p{0.22\linewidth}X@{}}
組織 & Qwen \\
種別 & - \\
時系列 & 2025-03-27 \\
\end{tabularx}
\smallskip
{\bfseries 一次資料から整理したtechnical points}
\begin{itemize}[leftmargin=1.5em,itemsep=0.35em]
\item \textbf{一次情報で確認できる事実}: 2025年3月27日、Qwenはtext・image・audio・video入力とstreaming speech出力を扱うQwen2.5-Omniを公開した。
\end{itemize}
\begin{samepage}
{\bfseries 一次資料}
\begingroup\sloppy
\Urlmuskip=0mu plus 2mu\relax
\begin{itemize}[leftmargin=1.5em,itemsep=0.25em]
\item {\scriptsize\url{https://qwenlm.github.io/blog/qwen2.5-omni/}}
\end{itemize}
\endgroup
\end{samepage}
\endgroup
\end{technicalnote}
\medskip

\begin{technicalnote}{Veo 3, Imagen 4, Flow, and Lyria 2}{主要資料}
\begingroup
% reader-facing Technical Notes break policy
\widowpenalty=10000
\clubpenalty=10000
\displaywidowpenalty=10000
\begin{tabularx}{\linewidth}{@{}>{\bfseries}p{0.22\linewidth}X@{}}
組織 & Google \\
種別 & - \\
時系列 & 2025-05-20 \\
\end{tabularx}
\smallskip
{\bfseries 一次資料から整理したtechnical points}
\begin{itemize}[leftmargin=1.5em,itemsep=0.35em]
\item \textbf{一次情報で確認できる事実}: 2025年5月20日、GoogleはVeo 3、Imagen 4、Flow、Lyria 2を含むgenerative media model/tool群を発表した。
\end{itemize}
\begin{samepage}
{\bfseries 一次資料}
\begingroup\sloppy
\Urlmuskip=0mu plus 2mu\relax
\begin{itemize}[leftmargin=1.5em,itemsep=0.25em]
\item {\scriptsize\url{https://blog.google/innovation-and-ai/products/generative-media-models-io-2025/}}
\end{itemize}
\endgroup
\end{samepage}
\endgroup
\end{technicalnote}
\medskip
